\section{Limitations, Ethics, And Transparency}

This journal draft is stronger than the preprint, but it is not yet a maximal claim paper.

\begin{itemize}[leftmargin=1.5em]
\item The strongest six-action positive result still comes from the bundled \texttt{freshness\_v1} benchmark rather than the public track.
\item The public FreshQA-derived benchmark is real and reproducible, but it is derived from weekly public releases and is therefore a public-track construction rather than a standard leaderboard split.
\item The public-track baselines named \texttt{self\_rag\_like}, \texttt{memllm\_like}, \texttt{wise\_like}, and \texttt{melo\_like} are approximate family baselines, not official reproductions.
\item The model matrix now includes two backbones rather than one, but it is still a small matrix.
\item Persistent memory and durable patches carry operational risks: stale evidence can be cached, unstable knowledge can be promoted, rollback can fail, and stored memories may create privacy or security concerns if deployed on sensitive data.
\end{itemize}

These limitations are exactly why the journal draft is useful in its current form. It moves the project from a narrow positive artifact package to a broader empirical study with explicit regime boundaries and reproducible next-step evidence.
